\input{../tex-inputs/latex-leseansicht-vorspann}

\section[Arthur Schnitzler an Gustav Schwarzkopf, 28. 4. 1922]{L04496 Arthur Schnitzler an Gustav Schwarzkopf, 28. 4. 1922}
\nopagebreak\mylabel{L04496v}
\rehead{ }\normalsize\beginnumbering\briefempfaengerindex{Schwarzkopf, Gustav@\textsc{Schwarzkopf, Gustav}!zzzSchnitzler, Arthur@\emph{von Arthur Schnitzler}!1922-04-281@{28. 4. 1922}|(be}
\toendnotes[C]{\smallbreak\pagebreak[2]}
\correspDesc{Versand  durch Arthur Schnitzler am 28. 4. 1922 in Den Haag
\newline{}Übermittlung  am 28. 4. 1922 in Den Haag
\newline{}Erhalt  durch Gustav Schwarzkopf im Zeitraum [29. 4. 1922
                  – 3. 5. 1922?] in Wien}\toendnotes[C]{\smallbreak}
\Standort{DLA, A:Schnitzler, HS.1985.1.1897.}
\physDesc{Bildpostkarte, 431 Zeichen
\newline{}Handschrift: Bleistift, deutsche Kurrent
\newline{}Versand: Stempel: »\nobreak{}\oindex{Den Haag@\textbf{Den Haag}|pwk}S' Gravenhage, 28. IV. 22, V, Gouda\oindex{Gouda@\textbf{Gouda}|pw}\nobreak{}«.  }\toendnotes[C]{\smallbreak}
\pstart
           \noindent{}{\pb}\textcolor{gray}{\textbf{’S GRAVENHAGE BLEIENBURG\oindex{Bleijenburg@\textbf{Bleijenburg}|pw}}}\pend
           \pstart{}{\pb}Hrn. \textsc{Gustav Schwarzkopf}\pend{}\pstart{}\textsc{Wien I}\oindex{I., Innere Stadt@\textbf{I., Innere Stadt}|pw}\pend{}\pstart{}\textsc{Tiefer Graben 17}\oindex{Wien@\textbf{Wien}!I., Innere Stadt@\textbf{I., Innere Stadt}!Tiefer Graben 17@\textbf{Tiefer Graben 17}|pw}\pend{}{\bigskip}\vspace{1em}
\pstart
           \raggedleft{}\textsc{Haag\oindex{Den Haag@\textbf{Den Haag}|pw}}{ }28. 4. 2\textcolor{gray}{2}\textcolor{gray}{}\pend
           \vspace{0.5em}
\pstart
           lieber Guſtav, ſeien Sie mir herzlich gegrüßt. \label{K_L04496-1v}\edtext{Zwei Vorleſungen\eventindex{Rotterdam@\textbf{Rotterdam}!Lesung von Lebendige Stunden, Weihnachts-Einkäufe, Das Schicksal des Freiherrn von Leisenbohg, 26.4.1922@Lesung von Lebendige Stunden, Weihnachts-Einkäufe, Das Schicksal des Freiherrn von Leisenbohg, 26.4.1922|pwv}\eventindex{Den Haag@\textbf{Den Haag}!Lesung von Die Hirtenflöte, Die letzten Masken, 27.4.1922@Lesung von Die Hirtenflöte, Die letzten Masken, 27.4.1922|pwv}}{\lemma{\textnormal{\emph{Zwei Vorlesungen}}}\Cendnote{\textnormal{Lesung von Lebendige Stunden, Weihnachts-Einkäufe, Das Schicksal des Freiherrn von Leisenbohg, 26.4.1922. In: A. S.: \emph{Kulturveranstaltungen}, \url{https://schnitzler-kultur.acdh.oeaw.ac.at/pmb183416.html} und Lesung von Die Hirtenflöte, Die letzten Masken, 27.4.1922. In: A. S.: \emph{Kulturveranstaltungen}, \url{https://schnitzler-kultur.acdh.oeaw.ac.at/pmb183422.html}.}}}\label{K_L04496-1}{ }ſind gut (und mit finanz.
               Erfolg-Garantie \textcolor{gray}{mich} überſchütten) vorbei, Rotterdam\oindex{Rotterdam@\textbf{Rotterdam}|pw}\eventindex{Rotterdam@\textbf{Rotterdam}!Lesung von Lebendige Stunden, Weihnachts-Einkäufe, Das Schicksal des Freiherrn von Leisenbohg, 26.4.1922@Lesung von Lebendige Stunden, Weihnachts-Einkäufe, Das Schicksal des Freiherrn von Leisenbohg, 26.4.1922|pwv}, Haag\oindex{Den Haag@\textbf{Den Haag}|pw}\eventindex{Den Haag@\textbf{Den Haag}!Lesung von Die Hirtenflöte, Die letzten Masken, 27.4.1922@Lesung von Die Hirtenflöte, Die letzten Masken, 27.4.1922|pwv}. Morgen{ }\label{K_L04496-2v}\edtext{Amſterdam\oindex{Amsterdam@\textbf{Amsterdam}|pw}\eventindex{Concertgebouw@\textbf{Concertgebouw}!Lesung von Der Mörder, Weihnachts-Einkäufe, Die letzten Masken, Excentric, 29.4.1922@Lesung von Der Mörder, Weihnachts-Einkäufe, Die letzten Masken, Excentric, 29.4.1922|pwv}}{\lemma{\textnormal{\emph{Amsterdam}}}\Cendnote{\textnormal{Lesung von Der Mörder, Weihnachts-Einkäufe, Die letzten Masken, Excentric, 29.4.1922. In: A. S.: \emph{Kulturveranstaltungen}, \url{https://schnitzler-kultur.acdh.oeaw.ac.at/pmb183430.html}.}}}\label{K_L04496-2}. »Einſamer Weg\pwindex{Schnitzler, Arthur 15. 5. 1862 Wien – 21. 10. 1931 ebd.@\textsc{Schnitzler, Arthur} (15. 5. 1862 Wien – 21. 10. 1931 ebd.)!einsame Weg. Schauspiel in fünf Akten@\strich\emph{Der einsame Weg. Schauspiel in fünf Akten}|pw}« – verſchoben – \label{K_L04496-3v}\edtext{beko{\geminationm}e nur – die Gefährtin\pwindex{Schnitzler, Arthur 15. 5. 1862 Wien – 21. 10. 1931 ebd.@\textsc{Schnitzler, Arthur} (15. 5. 1862 Wien – 21. 10. 1931 ebd.)!Gefährtin. Schauspiel in einem Akt@\strich\emph{Die Gefährtin. Schauspiel in einem Akt}|pw} zu{ }ſehen}{\lemma{\textnormal{\emph{bekomme … sehen}}}\Cendnote{\textnormal{Die Zeitung \emph{Vaderland}\pwindex{Het Vaderland@\emph{Het Vaderland}|pwk}
                  berichtete eine Woche zuvor über das Programm Schnitzlers auf seiner holländischen\oindex{Niederlande@\textbf{Niederlande}|pwk} Reise, unter anderem werde er am Nachmittag des
                     27. 4. 1922 der Generalprobe einer Inszenierung des \emph{Hamlet}\pwindex{Shakespeare, William 23.\,4.\,1564? Stratford-upon-Avon – 3.\,5.\,1616 ebd.@\textsc{Shakespeare, William} (23.\,4.\,1564? Stratford-upon-Avon – 3.\,5.\,1616 ebd.)!Hamlet@\strich\emph{Hamlet}|pwk} beiwohnen, bei der zuvor \emph{Die Gefährtin}\pwindex{Schnitzler, Arthur 15. 5. 1862 Wien – 21. 10. 1931 ebd.@\textsc{Schnitzler, Arthur} (15. 5. 1862 Wien – 21. 10. 1931 ebd.)!Gefährtin. Schauspiel in einem Akt@\strich\emph{Die Gefährtin. Schauspiel in einem Akt}|pwk} gegeben werde, \emph{Kunst en letteren. Arthur Schnitzler}\pwindex{Kunst en letteren. Arthur Schnitzler@\emph{Kunst en letteren. Arthur Schnitzler}|pwk}. In:
                        \emph{Het Vaderland}\pwindex{Het Vaderland@\emph{Het Vaderland}|pwk},
                        20. April 1922, Avondblad A, S. [5]. Dieser Probenbesuch
                  scheint nicht stattgefunden zu haben. Als Schnitzler am 30. 4. 1922 die Aufführung von
                     Hamlet\eventindex{Den Haag@\textbf{Den Haag}!Aufführung von Hamlet, 30.4.1922@Aufführung von Hamlet, 30.4.1922|pwk} besuchte, war nicht mehr die Rede von einem vorangestellten
                  Einakter in Form von \emph{Die Gefährtin}\pwindex{Schnitzler, Arthur 15. 5. 1862 Wien – 21. 10. 1931 ebd.@\textsc{Schnitzler, Arthur} (15. 5. 1862 Wien – 21. 10. 1931 ebd.)!Gefährtin. Schauspiel in einem Akt@\strich\emph{Die Gefährtin. Schauspiel in einem Akt}|pwk}. Regisseur
                     der \emph{Hamlet}\pwindex{Shakespeare, William 23.\,4.\,1564? Stratford-upon-Avon – 3.\,5.\,1616 ebd.@\textsc{Shakespeare, William} (23.\,4.\,1564? Stratford-upon-Avon – 3.\,5.\,1616 ebd.)!Hamlet@\strich\emph{Hamlet}|pwk}inszenierung\eventindex{Den Haag@\textbf{Den Haag}!Aufführung von Hamlet, 30.4.1922@Aufführung von Hamlet, 30.4.1922|pwkv} war Eduard Verkarde\pwindex{Verkade, Eduard Rutger 15.\,6.\,1878 Amsterdam – 11.\,2.\,1961 Breukelen@\textsc{Verkade, Eduard Rutger} (15.\,6.\,1878 Amsterdam – 11.\,2.\,1961 Breukelen)|pwk}, der daran
                     arbeitete, \emph{Der Einsamen Weg}\pwindex{Schnitzler, Arthur 15. 5. 1862 Wien – 21. 10. 1931 ebd.@\textsc{Schnitzler, Arthur} (15. 5. 1862 Wien – 21. 10. 1931 ebd.)!einsame Weg. Schauspiel in fünf Akten@\strich\emph{Der einsame Weg. Schauspiel in fünf Akten}|pwk} in den Niederlanden\oindex{Niederlande@\textbf{Niederlande}|pwk} auf die Bühne zu bringen. Möglicherweise hatte
                     er als Ersatz für die verschobene Inszenierung kurzfristig \emph{Die Gefährtin}\pwindex{Schnitzler, Arthur 15. 5. 1862 Wien – 21. 10. 1931 ebd.@\textsc{Schnitzler, Arthur} (15. 5. 1862 Wien – 21. 10. 1931 ebd.)!Gefährtin. Schauspiel in einem Akt@\strich\emph{Die Gefährtin. Schauspiel in einem Akt}|pwk} für Schnitzlers Besuch ins Programm nehmen wollen.}}}\label{K_L04496-3}. Die Gaſtfreundlichkeit hier
               über alle Begriffe und ohne jede Unbequemlichkeit.\pend
           
\pstart
           {\pb}Aus Kopenhagen\oindex{Kopenhagen@\textbf{Kopenhagen}|pw} hör
               ich nichts – fahre vielleicht danach nach Berlin\oindex{Berlin@\textbf{Berlin}|pw}.\pend
           
\pstart
           Ihr{\\[\baselineskip]}\spacefill\mbox{A.}\pend
           \leftskip=0em{}\selectlanguage{ngerman}\endnumbering\briefempfaengerindex{Schwarzkopf, Gustav@\textsc{Schwarzkopf, Gustav}!zzzSchnitzler, Arthur@\emph{von Arthur Schnitzler}!1922-04-281@{28. 4. 1922}|)be}\mylabel{L04496h}
\begin{anhang}
\end{anhang}\newcommand{\dateiname}{L04496}\newcommand{\titel}{Arthur Schnitzler an Gustav Schwarzkopf, 28. 4. 1922}\newcommand{\editorInnen}{Herausgegeben von Jahnke, SelmaMüller, Martin Anton}\input{../tex-inputs/latex-leseansicht-abspann}
